\documentclass[12pt]{article}
\usepackage[utf8]{inputenc}
\usepackage[T1]{fontenc}
\usepackage{hyperref}
\title{Matya Samo Language Resources for K-12 Curriculum Development}
\author{Compiled by Claude Sonnet 5.5 (claude-sonnet-5-5) for Andr\'as Kornai}
\date{}
\begin{document}
\maketitle

\section{Introduction}
Matya Samo (ISO 639-3 stj) is a Mande language of the Samo dialect cluster, spoken in Kossi and Sourou provinces of Burkina Faso (Nouna and Solenzo areas) and in Mali. Speaker estimates vary widely, from about 165,000 in one linguist's data note to over 240,000 in Joshua Project. Wikitongues rates it EGIDS 6a (vigorous). Search results turned up reference pages and a Scripture listing, but no school materials.

\section{Primary Educational Resources}
No primers, textbooks or school programmes were found in search results.

\section{Linguistic and Reference Materials}
\begin{itemize}
\item Samo language (Burkina), Wikipedia. \url{https://en.wikipedia.org/wiki/Samo_language_(Burkina)}
\item Matya Samo, Wikitongues language page. \url{https://wikitongues.org/languages/stj/}
\item Samo-Matya, numerals page from the Max Planck Institute for Evolutionary Anthropology (channumerals). \url{https://lingweb.eva.mpg.de/channumerals/Samo-Matya.htm}
\end{itemize}

\section{Religious and Community Resources}
\begin{itemize}
\item Joshua Project language page for stj. It reports that the New Testament is available and that an Android Bible app exists. \url{https://www.joshuaproject.net/index.php/languages/stj}
\item Joshua Project, Samo Matya in Burkina Faso. \url{https://www.joshuaproject.net/index.php/people_groups/14675/UV}
\item Joshua Project, Samo Matya in Mali. \url{https://www.joshuaproject.net/index.php/people_groups/14675/ML}
\end{itemize}

\section{Digital and Audio Resources}
No audio, video or corpus resources were found beyond the Android Bible app mentioned by Joshua Project. I did not find a link for the app.

\section{Curriculum Development Context}
Little educational infrastructure for Matya Samo appeared in the searches. A curriculum would need to build on the New Testament text and on community orthography work, neither of which I could verify here.

\end{document}